% ============================================================================
\section{Foundations: From Bits to Neurons}
\label{sec:foundations}
% ============================================================================

Before we can understand how to replace neural network arithmetic with
bitwise operations, we need a shared vocabulary. This section builds the
essential concepts from the ground up.

% ----------------------------------------------------------------------------
\subsection{What a CPU Actually Does Well}
\label{sec:cpu-ops}
% ----------------------------------------------------------------------------

Every processor---from a server CPU to a tiny microcontroller---has
a set of operations it can execute in a single clock cycle. The
cheapest and fastest of these are \concept{bitwise logic operations}:

\begin{figure}[htbp]
\centering
\begin{tikzpicture}[
  gate/.style={draw, minimum width=1.2cm, minimum height=0.8cm,
    font=\small\ttfamily\bfseries, rounded corners=2pt},
  bit/.style={font=\small\ttfamily},
  >=Stealth
]
  % AND gate
  \node[gate, fill=bitblue!15] (and) at (0,0) {AND};
  \node[bit, left=0.8cm] (a1) at ([yshift=0.2cm]and.west) {1};
  \node[bit, left=0.8cm] (a2) at ([yshift=-0.2cm]and.west) {0};
  \node[bit, right=0.8cm] (a3) at (and.east) {0};
  \draw[->] (a1) -- (and.west |- a1);
  \draw[->] (a2) -- (and.west |- a2);
  \draw[->] (and.east) -- (a3);
  \node[font=\scriptsize, below=0.1cm] at (and.south) {1 if both 1};

  % OR gate
  \node[gate, fill=bitgreen!15] (or) at (3.9,0) {OR};
  \node[bit, left=0.8cm] (b1) at ([yshift=0.2cm]or.west) {1};
  \node[bit, left=0.8cm] (b2) at ([yshift=-0.2cm]or.west) {0};
  \node[bit, right=0.8cm] (b3) at (or.east) {1};
  \draw[->] (b1) -- (or.west |- b1);
  \draw[->] (b2) -- (or.west |- b2);
  \draw[->] (or.east) -- (b3);
  \node[font=\scriptsize, below=0.1cm] at (or.south) {1 if either 1};

  % XOR gate
  \node[gate, fill=bitorange!15] (xor) at (7.8,0) {XOR};
  \node[bit, left=0.8cm] (c1) at ([yshift=0.2cm]xor.west) {1};
  \node[bit, left=0.8cm] (c2) at ([yshift=-0.2cm]xor.west) {0};
  \node[bit, right=0.8cm] (c3) at (xor.east) {1};
  \draw[->] (c1) -- (xor.west |- c1);
  \draw[->] (c2) -- (xor.west |- c2);
  \draw[->] (xor.east) -- (c3);
  \node[font=\scriptsize, below=0.1cm] at (xor.south) {1 if different};

  % NOT gate
  \node[gate, fill=bitpurple!15] (not) at (11.7,0) {NOT};
  \node[bit, left=0.8cm] (d1) at (not.west) {1};
  \node[bit, right=0.8cm] (d2) at (not.east) {0};
  \draw[->] (d1) -- (not.west);
  \draw[->] (not.east) -- (d2);
  \node[font=\scriptsize, below=0.1cm] at (not.south) {flip the bit};
\end{tikzpicture}
\caption{The four fundamental logic gates. Each operates on single bits
  and executes in one clock cycle. A 64-bit CPU applies these to 64~bits
  simultaneously.}
\label{fig:logic-gates}
\end{figure}

\begin{engineernote}
When a 64-bit CPU executes an AND instruction, it performs 64
independent AND operations in a single clock cycle. With SIMD
extensions (SSE, AVX-512, NEON), this scales to 128, 256, or even
512~bits at once. This massive parallelism is free---it is built
into the hardware.
\end{engineernote}

There is one more operation that matters enormously: \concept{popcount}
(population count). Given a string of bits, popcount counts how many
are~1. Most modern CPUs have a dedicated popcount instruction.

That instruction is useful, but it is not a universal performance shortcut.
On AVX2 systems, scalar popcount throughput and the absence of a native
vector-popcount instruction can make it a bottleneck in bitwise matrix
multiplication. A recent CPU study improved ternary and binary inference
by fusing packing with quantisation and replacing popcount-heavy kernels
with higher-throughput SIMD operations~\cite{fischer2025popcount}. Bitwise
models therefore need kernel-level measurement, not just operation counts.

\begin{figure}[htbp]
\centering
\begin{tikzpicture}[
  bit/.style={draw, minimum width=0.6cm, minimum height=0.6cm,
    font=\small\ttfamily, inner sep=0pt},
  >=Stealth
]
  % Bit string
  \foreach \b [count=\i from 0] in {1,0,1,1,0,1,0,1} {
    \ifnum\b=1
      \def\fillcol{bitblue!20}
    \else
      \def\fillcol{white}
    \fi
    \node[bit, fill=\fillcol]
      (b\i) at (\i*0.7, 0) {\b};
  }
  % Arrow
  \draw[->, thick] (3.5*0.7, -0.5) -- (3.5*0.7, -1.2);
  % Result
  \node[draw, rounded corners, fill=bitblue!10, minimum width=1.5cm,
    minimum height=0.7cm, font=\sffamily\bfseries] at (3.5*0.7, -1.7) {5};
  \node[font=\scriptsize\sffamily, right] at (5.5*0.7, -1.7)
    {(five 1-bits)};
  % Label
  \node[font=\sffamily\bfseries, above=0.3cm] at (3.5*0.7, 0.3)
    {popcount};
\end{tikzpicture}
\caption{The popcount operation counts the number of set bits (1s) in
  a binary word. Combined with XNOR, it becomes a similarity
  measure---see \cref{sec:bnn} for how this replaces multiplication
  in binary neural networks.}
\label{fig:popcount}
\end{figure}

Two more operations round out the bitwise toolkit:
\concept{bit shift} and \concept{bit rotation}. A \emph{left shift}
by $k$ positions moves every bit $k$ places toward the most
significant end, filling the vacated positions with zeros. A
\emph{right shift} does the reverse. Mathematically, shifting left by
$k$ is equivalent to multiplying by $2^k$, and shifting right by $k$
is equivalent to (integer) dividing by $2^k$---but the CPU does it in
a single cycle, without invoking the multiplier or divider circuit.

\begin{figure}[htbp]
\centering
\begin{tikzpicture}[
  bit/.style={draw, minimum width=0.6cm, minimum height=0.6cm,
    font=\small\ttfamily, inner sep=0pt},
  >=Stealth
]
  % Original
  \node[font=\sffamily\bfseries, left] at (-0.5, 0) {Original:};
  \foreach \b [count=\i from 0] in {0,0,1,0,1,1,0,1} {
    \ifnum\b=1
      \def\fillcol{bitblue!20}
    \else
      \def\fillcol{white}
    \fi
    \node[bit, fill=\fillcol] at (\i*0.7+2, 0) {\b};
  }
  \node[font=\scriptsize\sffamily, right] at (7.7, 0) {$= 45$};

  % Left shift by 2
  \node[font=\sffamily\bfseries, left] at (-0.5, -1.2) {$\ll 2$:};
  \foreach \b [count=\i from 0] in {1,0,1,1,0,1,0,0} {
    \ifnum\b=1
      \def\fillcol{bitorange!25}
    \else
      \def\fillcol{white}
    \fi
    \node[bit, fill=\fillcol] at (\i*0.7+2, -1.2) {\b};
  }
  \node[font=\scriptsize\sffamily, right] at (7.7, -1.2)
    {$= 180 = 45 \times 4$};

  % Right shift by 2
  \node[font=\sffamily\bfseries, left] at (-0.5, -2.4) {$\gg 2$:};
  \foreach \b [count=\i from 0] in {0,0,0,0,1,0,1,1} {
    \ifnum\b=1
      \def\fillcol{bitgreen!25}
    \else
      \def\fillcol{white}
    \fi
    \node[bit, fill=\fillcol] at (\i*0.7+2, -2.4) {\b};
  }
  \node[font=\scriptsize\sffamily, right] at (7.7, -2.4)
    {$= 11 = \lfloor 45 / 4 \rfloor$};
\end{tikzpicture}
\caption{Bit shifting. Left shift ($\ll$) multiplies by a power of~2;
  right shift ($\gg$) divides by a power of~2. Both execute in one
  clock cycle. Bit \emph{rotation} is similar but wraps bits around
  instead of discarding them.}
\label{fig:bitshift}
\end{figure}

\begin{engineernote}
Bit shifts turn multiplication and division by powers of~2 into
single-cycle operations. This is why they appear throughout the
integer-only methods in later sections: Shiftmax approximates the
exponential function using shifts (\cref{sec:softmax-problem}),
and integer normalisation replaces square roots and divisions with
iterative shift-and-add sequences
(\cref{sec:layernorm-problem}). Bit \emph{rotation} (wrapping
bits around instead of discarding them) is a standard CPU operation
used in cryptography and hashing but does not currently play a role
in any of the neural network paradigms discussed in this article.
\end{engineernote}

\subsection{The Cost of Arithmetic}
\label{sec:cost-arithmetic}

Not all CPU operations cost the same. \Cref{tab:op-cost} shows
approximate energy costs for different operations in a modern 45\,nm
chip~\cite{horowitz2014energy}.

\begin{table}[htbp]
\centering
\caption{Energy cost of arithmetic operations at a 45\,nm process
  node~\cite{horowitz2014energy}. Modern processors use much smaller
  nodes (5--7\,nm), which lowers \emph{absolute} energies roughly
  5--10$\times$~\cite{stillmaker2017scaling}. However, the \emph{ratios}
  between operation types are determined by algorithmic complexity (a
  multiplier uses $O(n^2)$ gates, an adder $O(n)$, a bitwise operation
  $O(1)$ per bit) and therefore remain approximately constant across
  process generations.}
\label{tab:op-cost}
\begin{tabular}{@{}lrr@{}}
  \toprule
  \textbf{Operation}      & \textbf{Energy (pJ)} & \textbf{Relative Cost} \\
  \midrule
  Bitwise AND/OR/XOR      & 0.1                  & 1$\times$ \\
  Bit shift                & 0.1                  & 1$\times$ \\
  Integer add (8-bit)      & 0.03                 & 0.3$\times$ \\
  Integer multiply (8-bit) & 0.2                  & 2$\times$ \\
  Integer add (32-bit)     & 0.1                  & 1$\times$ \\
  Integer multiply (32-bit)& 3.1                  & 31$\times$ \\
  FP add (32-bit)          & 0.9                  & 9$\times$ \\
  FP multiply (32-bit)     & 3.7                  & 37$\times$ \\
  Memory access (SRAM)     & 5.0                  & 50$\times$ \\
  Memory access (DRAM)     & 640.0                & 6400$\times$ \\
  \bottomrule
\end{tabular}
\end{table}

\begin{keyinsight}
The most expensive thing a chip does is not computation---it is
\emph{moving data}. A DRAM memory access costs 6400$\times$ more
energy than a logic gate. This is the ``memory wall,'' and it means
that any approach which keeps data close to the logic (small models,
local lookups) wins enormously on energy, even if it does more
operations.
\end{keyinsight}

% ----------------------------------------------------------------------------
\subsection{The Precision Spectrum: From 1-Bit to 8-Bit}
\label{sec:precision-spectrum}
% ----------------------------------------------------------------------------

The paradigms in this article differ not only in \emph{which}
operations they use, but in how many bits represent each value.
This distinction is important because it determines what bitwise
operations can meaningfully compute.

\subsubsection{Why Most Approaches Booleanise to 1~Bit}

If every input and every weight is a single bit, then AND, OR, XOR,
and NOT are the \emph{natural} operations---they map exactly onto
the logic of combining two Boolean propositions. XNOR on two 1-bit
values tells you whether they agree. Popcount on a vector of XNOR
results tells you how many positions agree---a Hamming
similarity.\footnote{The \emph{Hamming distance} between two binary
strings is the number of positions at which the bits differ. It was
introduced by Richard Hamming in~1950 for error-detecting codes.
\emph{Hamming similarity} is the complement: the number of positions
that \emph{agree}. For two $n$-bit strings $a$ and~$b$:
$\text{similarity} = n - \text{distance}
= \operatorname{popcount}(\operatorname{XNOR}(a,b))$.}
This is maximally hardware-efficient: a 64-bit CPU processes 64
independent 1-bit operations in a single instruction.

The price is \concept{precision loss}. For MNIST, collapsing each
pixel to 1~bit (ink or background) loses surprisingly little
information---the images are nearly binary already. But for richer
inputs (colour images, audio spectrograms, word embeddings),
reducing to 1~bit per feature is destructive.

\subsubsection{The Positional Encoding Problem}

Why not simply apply XNOR to 8-bit values? You can---but the
result is not what you might expect. An 8-bit integer stores its
value in \emph{positional binary encoding}: bit~7 is worth~128,
bit~0 is worth~1. Bitwise operations treat all bit positions
identically---they do not know that bit~7 matters 128$\times$ more
than bit~0.

\begin{analogy}
Applying XOR to two 8-bit integers is like comparing two
five-digit numbers by checking how many \emph{digits} match,
regardless of position. The numbers 50{,}001 and 50{,}002 differ
in one digit (the units), but so do 50{,}001 and 10{,}001 (the
ten-thousands). Position-blind comparison treats these as equally
different, even though the first pair is numerically close and the
second pair is far apart.
\end{analogy}

To make bitwise operations numerically meaningful on multi-bit
values, you need one of three approaches:

\begin{enumerate}[leftmargin=2em]
  \item \textbf{Carry propagation}---let the result of one bit
    position influence the next. This is exactly what integer
    addition does: it is implemented as a chain of XOR gates (for
    the sum bits) and AND gates (for the carries). In other words,
    integer addition \emph{is} a structured sequence of bitwise
    operations with inter-bit dependencies.

  \item \textbf{Unary (thermometer) encoding}---represent the
    value~5 not as \texttt{0101} but as \texttt{11111000} (five~1s
    followed by zeros). Now bitwise AND gives the minimum of two
    values, OR gives the maximum, and Hamming distance correlates
    with numerical distance. The cost: a $k$-bit value becomes a
    $(2^k - 1)$-bit thermometer, a large expansion.

  \item \textbf{Accept the approximation}---use Hamming distance on
    raw binary as a rough similarity measure. This works
    surprisingly well when combined with large vectors and
    averaging (high-dimensional statistics smooth out the
    positional noise).
\end{enumerate}

\subsubsection{The Precision--Operation Map}

\Cref{tab:precision-map} shows how different approaches in this
article combine data precision with operation types. The field has
split into two camps: \emph{pure logic} (1-bit data, logic gates)
and \emph{integer quantisation} (multi-bit data, integer arithmetic).

\begin{table}[htbp]
\centering
\caption{The precision spectrum. Approaches differ in how many bits
  represent data and weights, which determines what operations are
  efficient.}
\label{tab:precision-map}
\small
\begin{tabular}{@{}p{2.6cm}ccp{3.7cm}l@{}}
  \toprule
  \textbf{Approach} & \textbf{Data} & \textbf{Weights}
    & \textbf{Core operations} & \textbf{Used by} \\
  \midrule
  Full binary
    & 1 bit & 1 bit
    & AND, OR, XOR, NOT
    & BNN, TM, LGN \\[3pt]
  Thermometer
    & $k$ bits (unary) & 1 bit
    & Logic on expanded features
    & Advanced TM \\[3pt]
  Ternary weights
    & 8 bit & 1.58 bit
    & Integer add/subtract
    & BitNet b1.58 \\[3pt]
  Low-bit integer
    & 4--8 bit & 4--8 bit
    & Integer + bit shifts
    & I-BERT, I-ViT \\[3pt]
  LUT deployment
    & any & (compiled)
    & Memory read
    & KAN/LUT, ULEEN \\
  \bottomrule
\end{tabular}
\end{table}

\begin{keyinsight}
Notice the gap in \cref{tab:precision-map}: \textbf{no approach
combines multi-bit data with pure bitwise operations (no integer
addition).} The ``pure logic'' camp booleanises everything to 1~bit;
the ``integer'' camp keeps multi-bit data but relies on integer
addition with carry propagation. The middle ground---multi-bit data
processed only with shifts, masks, and position-independent logic---is
largely unexplored. Whether novel encodings (thermometer, Gray code,
or something entirely new) could make this viable is an open research
question. Classical integer addition costs $9\times$ the energy of a
bitwise operation (\cref{tab:op-cost}); if it could be avoided while
retaining multi-bit precision, the efficiency gain would be
significant.
\end{keyinsight}

% ----------------------------------------------------------------------------
\subsection{How a Traditional Neuron Works}
\label{sec:traditional-neuron}
% ----------------------------------------------------------------------------

A single artificial neuron does three things:

\begin{enumerate}[leftmargin=2em]
  \item \textbf{Multiply} each input $x_i$ by a learned weight $w_i$.
  \item \textbf{Accumulate} (sum) all the products and add a bias $b$.
  \item \textbf{Activate} by passing the sum through a nonlinear function $\sigma$.
\end{enumerate}

\begin{equation}
  y = \sigma\!\left(\sum_{i=1}^{n} w_i \cdot x_i + b\right)
  \label{eq:neuron}
\end{equation}

\begin{figure}[htbp]
\centering
\begin{tikzpicture}[
  input/.style={circle, draw, minimum size=0.7cm, font=\small},
  neuron/.style={circle, draw, minimum size=1.2cm, font=\small,
    fill=bitblue!10, line width=1pt},
  >=Stealth
]
  % Inputs
  \foreach \i/\v in {1/x_1, 2/x_2, 3/x_3} {
    \node[input] (in\i) at (0, 2-\i) {$\v$};
  }
  \node at (0, -1.5) {$\vdots$};
  \node[input] (inn) at (0, -2.5) {$x_n$};

  % Neuron
  \node[neuron] (n) at (4, -0.25) {$\Sigma$};

  % Activation
  \node[draw, rounded corners, minimum width=1cm, minimum height=0.8cm,
    fill=bitorange!10, font=\small] (act) at (6.5, -0.25) {$\sigma$};

  % Output
  \node[right=0.8cm, font=\small] (out) at (act.east) {$y$};

  % Connections with weights
  \foreach \i in {1,2,3} {
    \draw[->, thick] (in\i) -- node[above, font=\scriptsize, sloped]
      {$w_\i$} (n);
  }
  \draw[->, thick] (inn) -- node[above, font=\scriptsize, sloped]
    {$w_n$} (n);

  % Neuron to activation
  \draw[->, thick] (n) -- (act);
  \draw[->, thick] (act) -- (out);

  % Bias
  \draw[->, thick, dashed] (4, -2.5) -- node[right, font=\scriptsize]
    {bias $b$} (n);

  % Labels
  \node[font=\scriptsize\sffamily, below=0.1cm, bitgray] at (2, -3.2)
    {multiply-accumulate (MAC)};
  \draw[decorate, decoration={brace, amplitude=5pt, mirror},
    bitgray, thick] (0.8, -3) -- (5.2, -3);
\end{tikzpicture}
\caption{A traditional artificial neuron. The multiply-accumulate step
  dominates both computation time and energy cost. Every paradigm in
  this article replaces this step with something cheaper.}
\label{fig:traditional-neuron}
\end{figure}

\begin{analogy}
Think of a neuron as a judge at a talent show. Each judge gives a
score (the input $x_i$), and each judge has a reputation weight
($w_i$)---a famous expert's opinion counts more. The final score is
the weighted average. The activation function is the threshold: above
a certain score, the act goes through; below, it is rejected.

The MAC approach computes this weighted average with precise
floating-point arithmetic. The bitwise approaches we will explore
are like replacing the nuanced scoring with a simpler system:
thumbs-up or thumbs-down from each judge, then just counting votes.
\end{analogy}

% ----------------------------------------------------------------------------
\subsection{What MNIST Looks Like}
\label{sec:mnist}
% ----------------------------------------------------------------------------

The MNIST dataset~\cite{lecun1998mnist} is the standard first benchmark
for any new neural network approach. It consists of 70,000 handwritten
digit images (60,000 for training, 10,000 for testing), each
28$\times$28 pixels in grayscale (values 0--255).

\begin{figure}[htbp]
\centering
\begin{tikzpicture}[scale=0.16]
  % Draw a stylised "3" as a pixel grid
  \fill[bitblue!5] (0,0) rectangle (28,28);
  \draw[bitgray!30, very thin, step=1] (0,0) grid (28,28);

  % Approximate pixels for digit "3"
  \foreach \x/\y/\c in {
    10/24/80, 11/24/200, 12/24/255, 13/24/255, 14/24/255, 15/24/200, 16/24/80,
    16/23/200, 17/23/255, 18/23/200,
    17/22/200, 18/22/255,
    16/21/200, 17/21/255,
    14/20/200, 15/20/255, 16/20/255,
    14/19/255, 15/19/200,
    16/18/200, 17/18/255, 18/18/200,
    18/17/255, 19/17/200,
    18/16/255, 19/16/200,
    17/15/200, 18/15/255,
    15/14/200, 16/14/255, 17/14/255,
    11/13/80, 12/13/200, 13/13/255, 14/13/255, 15/13/200
  }{
    \pgfmathsetmacro{\intensity}{\c/255*100}
    \fill[black!\intensity] (\x,\y) rectangle ++(1,1);
  }

  % Arrow
  \draw[-{Stealth}, very thick, bitgray] (30, 14) -- (34, 14);

  % Vector representation
  \node[font=\small\sffamily, anchor=south] at (37, 23) {784 inputs};
  \draw[thick] (35, 5) -- (35, 22);
  \foreach \y/\v in {21/0.00, 19.5/0.00, 18/0.31, 16.5/0.78, 15/1.00,
    13.5/0.78, 12/\ldots, 10.5/0.00, 9/0.00, 7.5/\ldots, 6/0.85} {
    \node[font=\tiny\ttfamily, right] at (36, \y) {\v};
  }

  % Arrow to output
  \draw[-{Stealth}, very thick, bitgray] (41, 14) -- (45, 14);

  % Output
  \node[font=\small\sffamily, anchor=south] at (51, 23) {Output};
  \foreach \d [count=\i from 0] in {0,1,2,3,4,5,6,7,8,9} {
    \pgfmathsetmacro{\ypos}{22 - \i * 2.0}
    \node[font=\small\ttfamily, right] at (46, \ypos) {\d};
    \ifnum\d=3
      \fill[bitgreen!60] (49, {\ypos-0.5}) rectangle (53, {\ypos+0.5});
      \node[font=\tiny\sffamily, right] at (53.5, \ypos) {0.97};
    \else
      \fill[bitgray!15] (49, {\ypos-0.5}) rectangle (49.8, {\ypos+0.5});
    \fi
  }
\end{tikzpicture}
\caption{The MNIST task: a 28$\times$28 grayscale image is flattened
  into 784 values (normalised to $[0,1]$), and the network must output
  which digit (0--9) it represents. The standard floating-point
  baseline achieves 98--99\% accuracy.}
\label{fig:mnist}
\end{figure}

For bitwise networks, the first question is how to represent these
pixel values as bits. The simplest approach is \concept{thresholding}:
if a pixel is above 127, it becomes a~1; otherwise a~0. More
sophisticated approaches use \concept{thermometer encoding} (multiple
bits per pixel) or learned binarisation~\cite{lab2023}. We will revisit this in each
paradigm's section.

% ----------------------------------------------------------------------------
\subsection{What ``Training'' Means}
\label{sec:training-overview}
% ----------------------------------------------------------------------------

Training a neural network means adjusting its parameters (weights,
thresholds, gate types, spline coefficients---depending on the paradigm)
so that the network's outputs match the \emph{known correct answers} in
a labelled training set. This is \concept{supervised learning}: each
training example comes with the right answer (e.g.\ an image of a~3
paired with the label ``3''), and the network learns by comparing its
prediction to that label. In traditional networks, this comparison-and-correction
cycle uses \concept{gradient descent}:

\begin{enumerate}[leftmargin=2em]
  \item Run inputs through the network (\emph{forward pass}).
  \item Measure how wrong the output is by comparing it to the known
    correct answer---this score is called the \emph{loss}.
  \item Compute how each parameter contributed to the error
    (\emph{backward pass / backpropagation}).
  \item Nudge each parameter a tiny step in the direction that
    reduces the error.
\end{enumerate}

\begin{engineernote}
Backpropagation works by applying the chain rule of calculus: starting
from the loss, it computes the derivative of each layer's output with
respect to that layer's parameters, then chains these derivatives
together to propagate the error signal all the way back to the first
layer. This requires that every operation in the network be
\emph{differentiable}---that we can compute how a tiny change in its
input affects its output. Logic gates (AND, OR, XOR) are \emph{not}
differentiable: their output jumps discretely from~0 to~1 with no
smooth transition, so the derivative is either zero or undefined. This
is the central challenge of bitwise neural networks, and each paradigm
solves it differently:
\begin{itemize}[leftmargin=1em]
  \item BNNs use the ``straight-through estimator''~\cite{bengio2013ste}---pretend the
    derivative exists (\cref{sec:bnn}).
  \item Tsetlin Machines~\cite{granmo2018tsetlin} avoid gradients entirely, using
    \emph{reinforcement learning}---a training paradigm where each
    component receives simple reward/penalty signals (``that clause
    was useful'' or ``that clause was not'') and adjusts its behaviour
    accordingly, with no derivatives required (\cref{sec:tsetlin}).
  \item LGNs~\cite{petersen2022difflogic} use a continuous relaxation that converges to discrete
    gates (\cref{sec:lgn}).
  \item KANs train in continuous space and discretise only at
    deployment (\cref{sec:kan}).
\end{itemize}
\end{engineernote}

With this foundation, we are ready to explore each paradigm in detail.
